\documentclass[11pt,a4paper]{article}
\usepackage[utf8]{inputenc}
\usepackage[T1]{fontenc}
\usepackage{amsmath, amssymb}
\usepackage{amsthm}
\usepackage[margin=2.5cm]{geometry}
\usepackage{hyperref}
\hypersetup{pdftitle={Energy Decay in a Damped Harmonic Oscillator}, pdfauthor={Lena Ortiz}}

\newtheorem{theorem}{Theorem}

\title{Energy Decay in a Damped Harmonic Oscillator}
\author{Lena Ortiz\\ \small Department of Physics, Northfield College}
\date{October 2026}

\begin{document}
\maketitle

\begin{abstract}
We study the free motion of a mass on a linear spring subject to viscous damping. We derive the closed-form solution in the underdamped regime, show that the mechanical energy decays monotonically, and compare the analytic solution with two simple numerical schemes.
\end{abstract}

\section{Introduction}\label{sec:intro}

The damped harmonic oscillator is the simplest model of a system that loses energy to its surroundings~\cite{ref1}. Its equation of motion is
\begin{equation}
  m\ddot{x}+c\dot{x}+kx=0,
  \label{eq:gauss}
\end{equation}
where $m>0$ is the mass, $c\geq0$ the damping coefficient and $k>0$ the spring constant.\footnote{All quantities are in SI units.} Dividing by $m$ and writing $\omega_{0}=\sqrt{k/m}$ and $\gamma=c/(2m)$ gives the \emph{standard form}
\begin{equation}
  \ddot{x}+2\gamma\dot{x}+\omega_{0}^{2}x=0.
  \label{eq:2}
\end{equation}

\section{Solution}

Substituting $x=e^{\lambda t}$ into~\eqref{eq:2} gives the characteristic equation $\lambda^{2}+2\gamma\lambda+\omega_{0}^{2}=0$, whose roots are
\begin{equation}
  \lambda_{\pm}=-\gamma\pm\sqrt{\gamma^{2}-\omega_{0}^{2}}.
\end{equation}
The behaviour depends on the sign of the discriminant:
\begin{equation*}
  regime=\begin{cases}
 & \gamma<\omega_{0}, \\
 & \gamma=\omega_{0},
\end{cases}
\end{equation*}
In the underdamped case let $\omega_{d}=\sqrt{\omega_{0}^{2}-\gamma^{2}}$. With $x(0)=x_{0}$ and $\dot{x}(0)=0$ the solution is
\begin{equation}
  x(t)=x_{0}e^{-\gamma t}(\cos)\omega_{d}t+\frac{\gamma}{\omega_{d}\sin}\omega_{d}t.
\end{equation}
Equivalently, with $\mathbf{}y=(x,\dot{x})^{T}$, equation~\eqref{eq:2} becomes the first-order system
\begin{equation}
  \dot{y}=\begin{pmatrix}
0 & 1 \\
-\omega_{0}^{2}-2\gamma & 
\end{pmatrix}y.
\end{equation}

\begin{theorem}
  The energy $E(t)=\frac{1}{2}m\dot{x}^{2}+\frac{1}{2}kx^{2}$ never increases.
\end{theorem}

\begin{thebibliography}{9}
\bibitem{ref1} H. Goldstein, C. Poole and J. Safko, \emph{Classical Mechanics, 3rd ed.}, Addison-Wesley, 2002.
\bibitem{ref2} E. Hairer, S. P. Nørsett and G. Wanner, \emph{Solving Ordinary Differential Equations I, 2nd ed.}, Springer, 1993.
\end{thebibliography}

\end{document}
